% ============================================================
%  SAMURAI · МПС, объяснённая с нуля
%  Общий preamble (пакеты, окружения, макросы).
%
%  Файл подключается из main.tex и из «mini-main» каждой главы,
%  чтобы любую главу можно было собрать в отдельный PDF.
% ============================================================

% ── Кодировка и язык ─────────────────────────────────────────────
\usepackage[utf8]{inputenc}
\usepackage[T2A]{fontenc}
\usepackage[russian, english]{babel}

% ── Геометрия страницы ───────────────────────────────────────────
\usepackage[
  top=25mm, bottom=25mm,
  left=30mm, right=20mm
]{geometry}

% ── Математика ───────────────────────────────────────────────────
\usepackage{amsmath, amssymb, amsthm}
\usepackage{mathtools}
\usepackage{bm}

% ── Графика и цвет ───────────────────────────────────────────────
\usepackage{xcolor}
\usepackage{graphicx}
\usepackage{tikz}
\usetikzlibrary{arrows.meta, positioning, shapes.geometric, calc, fit,
                decorations.pathmorphing, decorations.markings, backgrounds}

% ── Код ──────────────────────────────────────────────────────────
\usepackage{listings}
\usepackage{courier}

\definecolor{codebg}{RGB}{245,247,250}
\definecolor{codeframe}{RGB}{180,190,210}
\definecolor{keyword}{RGB}{0,100,200}
\definecolor{comment}{RGB}{100,130,100}
\definecolor{string}{RGB}{163,21,21}
\definecolor{number}{RGB}{9,134,88}

% Поддержка кириллицы в listings.
\lstset{
  extendedchars=true,
  inputencoding=utf8,
  literate=
    {а}{{\selectfont\char224}}1{б}{{\selectfont\char225}}1
    {в}{{\selectfont\char226}}1{г}{{\selectfont\char227}}1
    {д}{{\selectfont\char228}}1{е}{{\selectfont\char229}}1
    {ё}{{\selectfont\char184}}1{ж}{{\selectfont\char230}}1
    {з}{{\selectfont\char231}}1{и}{{\selectfont\char232}}1
    {й}{{\selectfont\char233}}1{к}{{\selectfont\char234}}1
    {л}{{\selectfont\char235}}1{м}{{\selectfont\char236}}1
    {н}{{\selectfont\char237}}1{о}{{\selectfont\char238}}1
    {п}{{\selectfont\char239}}1{р}{{\selectfont\char240}}1
    {с}{{\selectfont\char241}}1{т}{{\selectfont\char242}}1
    {у}{{\selectfont\char243}}1{ф}{{\selectfont\char244}}1
    {х}{{\selectfont\char245}}1{ц}{{\selectfont\char246}}1
    {ч}{{\selectfont\char247}}1{ш}{{\selectfont\char248}}1
    {щ}{{\selectfont\char249}}1{ъ}{{\selectfont\char250}}1
    {ы}{{\selectfont\char251}}1{ь}{{\selectfont\char252}}1
    {э}{{\selectfont\char253}}1{ю}{{\selectfont\char254}}1
    {я}{{\selectfont\char255}}1
    {А}{{\selectfont\char192}}1{Б}{{\selectfont\char193}}1
    {В}{{\selectfont\char194}}1{Г}{{\selectfont\char195}}1
    {Д}{{\selectfont\char196}}1{Е}{{\selectfont\char197}}1
    {Ё}{{\selectfont\char168}}1{Ж}{{\selectfont\char198}}1
    {З}{{\selectfont\char199}}1{И}{{\selectfont\char200}}1
    {Й}{{\selectfont\char201}}1{К}{{\selectfont\char202}}1
    {Л}{{\selectfont\char203}}1{М}{{\selectfont\char204}}1
    {Н}{{\selectfont\char205}}1{О}{{\selectfont\char206}}1
    {П}{{\selectfont\char207}}1{Р}{{\selectfont\char208}}1
    {С}{{\selectfont\char209}}1{Т}{{\selectfont\char210}}1
    {У}{{\selectfont\char211}}1{Ф}{{\selectfont\char212}}1
    {Х}{{\selectfont\char213}}1{Ц}{{\selectfont\char214}}1
    {Ч}{{\selectfont\char215}}1{Ш}{{\selectfont\char216}}1
    {Щ}{{\selectfont\char217}}1{Ъ}{{\selectfont\char218}}1
    {Ы}{{\selectfont\char219}}1{Ь}{{\selectfont\char220}}1
    {Э}{{\selectfont\char221}}1{Ю}{{\selectfont\char222}}1
    {Я}{{\selectfont\char223}}1
    {—}{{\textemdash}}1{–}{{\textendash}}1{«}{{\guillemotleft}}1
    {»}{{\guillemotright}}1{№}{{N\textsuperscript{o}}}1
    {²}{{\textsuperscript{2}}}1{·}{{$\cdot$}}1{°}{{$^\circ$}}1
    {≤}{{$\leq$}}1{≥}{{$\geq$}}1{≠}{{$\neq$}}1{≈}{{$\approx$}}1
    {±}{{$\pm$}}1{×}{{$\times$}}1{÷}{{$\div$}}1{√}{{$\surd$}}1
    {→}{{$\to$}}1{←}{{$\leftarrow$}}1{⇒}{{$\Rightarrow$}}1
    {⇔}{{$\Leftrightarrow$}}1
    {θ}{{$\theta$}}1{φ}{{$\varphi$}}1{ω}{{$\omega$}}1
    {Φ}{{$\Phi$}}1{Γ}{{$\Gamma$}}1{Σ}{{$\Sigma$}}1{Ω}{{$\Omega$}}1
    {ε}{{$\varepsilon$}}1{σ}{{$\sigma$}}1{μ}{{$\mu$}}1{ν}{{$\nu$}}1
    {ρ}{{$\rho$}}1{δ}{{$\delta$}}1{η}{{$\eta$}}1{ψ}{{$\psi$}}1
    {α}{{$\alpha$}}1{β}{{$\beta$}}1{γ}{{$\gamma$}}1
    {τ}{{$\tau$}}1{π}{{$\pi$}}1{Δ}{{$\Delta$}}1{λ}{{$\lambda$}}1
}

\lstdefinestyle{pythonstyle}{
  language=Python,
  backgroundcolor=\color{codebg},
  frame=single, rulecolor=\color{codeframe}, framesep=4pt,
  basicstyle=\footnotesize\ttfamily,
  keywordstyle=\color{keyword}\bfseries,
  commentstyle=\color{comment}\itshape,
  stringstyle=\color{string},
  numberstyle=\scriptsize\ttfamily\color{gray},
  numbers=left, numbersep=10pt, xleftmargin=2.5em,
  showstringspaces=false, breaklines=true,
  tabsize=4, captionpos=b,
  morekeywords={True, False, None, self},
}

\lstdefinestyle{matlabstyle}{
  language=Matlab,
  backgroundcolor=\color{codebg},
  frame=single, rulecolor=\color{codeframe}, framesep=4pt,
  basicstyle=\footnotesize\ttfamily,
  keywordstyle=\color{keyword}\bfseries,
  commentstyle=\color{comment}\itshape,
  stringstyle=\color{string},
  numberstyle=\scriptsize\ttfamily\color{gray},
  numbers=left, numbersep=10pt, xleftmargin=2.5em,
  showstringspaces=false, breaklines=true,
  tabsize=2, captionpos=b,
}

\lstdefinestyle{tsstyle}{
  language=Java,                % близкий к TS лексер
  backgroundcolor=\color{codebg},
  frame=single, rulecolor=\color{codeframe}, framesep=4pt,
  basicstyle=\footnotesize\ttfamily,
  keywordstyle=\color{keyword}\bfseries,
  commentstyle=\color{comment}\itshape,
  stringstyle=\color{string},
  numberstyle=\scriptsize\ttfamily\color{gray},
  numbers=left, numbersep=10pt, xleftmargin=2.5em,
  showstringspaces=false, breaklines=true,
  tabsize=2, captionpos=b,
  morekeywords={const, let, function, export, import, from, type,
                interface, number, string, boolean, return, if, else},
}

\lstset{style=pythonstyle}

% ── Таблицы ──────────────────────────────────────────────────────
\usepackage{booktabs}
\usepackage{array}
\usepackage{longtable}
\usepackage{float}

% ── Гиперссылки ──────────────────────────────────────────────────
\usepackage[
  colorlinks=true,
  linkcolor=blue!70!black,
  citecolor=green!60!black,
  urlcolor=blue!80!black,
  bookmarks=true,
  bookmarksnumbered=true
]{hyperref}

% ── Колонтитулы ──────────────────────────────────────────────────
\usepackage{fancyhdr}
\setlength{\headheight}{14pt}     % fancyhdr требует >= 13.6pt для русского fontenc
\addtolength{\topmargin}{-2pt}
\pagestyle{fancy}
\fancyhf{}
\fancyhead[LE,RO]{\small\thepage}
\fancyhead[LO]{\small\nouppercase{\rightmark}}
\fancyhead[RE]{\small\nouppercase{\leftmark}}
\renewcommand{\headrulewidth}{0.4pt}

% ── Параграфы ────────────────────────────────────────────────────
\usepackage{parskip}
\setlength{\parindent}{0pt}

% ── Окружения ────────────────────────────────────────────────────
\usepackage{tcolorbox}
\tcbuselibrary{skins, breakable}

% «Идея»: интуиция/мотивация перед формальной частью.
\newtcolorbox{ideabox}[1][]{
  enhanced, breakable,
  colback=blue!4!white, colframe=blue!55!black,
  fonttitle=\bfseries, title=Идея, #1
}
% «Формула»: ключевое уравнение, на которое потом ссылаемся.
\newtcolorbox{formulabox}[1][]{
  enhanced, breakable,
  colback=blue!3!white, colframe=blue!50!black,
  fonttitle=\bfseries, title=Формула, #1
}
% «Примечание»: оранжевая «полочка», объяснение в сторону.
\newtcolorbox{notebox}[1][]{
  enhanced,
  colback=orange!8!white, colframe=orange!70!black,
  fonttitle=\bfseries, title=Примечание, #1
}
% «Внимание»: красная коробка — типичная ошибка / подводный камень.
\newtcolorbox{warnbox}[1][]{
  enhanced,
  colback=red!4!white, colframe=red!70!black,
  fonttitle=\bfseries, title=Внимание, #1
}
% «В коде»: серая «полочка» — ссылка на файл/строку с пояснением.
\newtcolorbox{codebox}[1][]{
  enhanced, breakable,
  colback=gray!6!white, colframe=gray!50!black,
  fonttitle=\bfseries\ttfamily, title=В коде, #1
}
% «Пример»: численный пример (зелёная).
\newtcolorbox{examplebox}[1][]{
  enhanced, breakable,
  colback=green!4!white, colframe=green!45!black,
  fonttitle=\bfseries, title=Пример, #1
}
% «Что если?»: фиолетовая — учебное «а что будет, если...».
\newtcolorbox{whatifbox}[1][]{
  enhanced, breakable,
  colback=violet!5!white, colframe=violet!55!black,
  fonttitle=\bfseries, title=Что если?, #1
}

% ── Нумерация уравнений по главам ────────────────────────────────
\numberwithin{equation}{chapter}

% ── Макросы ──────────────────────────────────────────────────────
\newcommand{\filepath}[1]{\texttt{\small\color{gray!80!black}#1}}
\newcommand{\vect}[1]{\bm{#1}}
\newcommand{\mat}[1]{\mathbf{#1}}
\newcommand{\T}{^{\mathsf{T}}}
\newcommand{\Real}{\mathbb{R}}
\newcommand{\diag}{\operatorname{diag}}
\newcommand{\rank}{\operatorname{rank}}
\newcommand{\sign}{\operatorname{sign}}
\DeclareMathOperator*{\argmin}{arg\,min}
\DeclareMathOperator*{\argmax}{arg\,max}
% \half — короткое 1/2.
\newcommand{\half}{\tfrac{1}{2}}
